% !PRINTED 569
\weberpage{569}
%<<CONT
setzen, worin $\psi$ ein ganzes Functional, $\varepsilon$ eine Einheit ist, und
erhalten aus (14)
\weq{15}{N_a(\psi) < \sqrt{\pm 0{,}404\, \Delta}.}

Da wir f\"ur $\varphi$ jedes ganze Functional setzen k\"onnen, so kann
$\psi$ ein Repr\"asentant eines Ideals einer beliebigen Classe sein, und
wir haben damit den wichtigen Satz bewiesen:

3. \gesperrt{In jeder Idealclasse k\"onnen wir einen Repr\"asen-}
\gesperrt{tanten finden, dessen Norm kleiner als die}
\gesperrt{Quadratwurzel aus dem absoluten Werthe der}
\gesperrt{Grundzahl des K\"orpers ist.}

Die im \S.\,153 zum Beweise der Endlichkeit der Classenzahl
benutzte obere Grenze f\"ur die Norm eines Ideals einer Classe
wird hierdurch in viel sch\"arferer Weise bestimmt.

Aber noch einen anderen sehr bemerkenswerthen Schluss
gestattet die Ungleichung (15). Da n\"amlich $N_a(\psi)$ eine positive
ganze Zahl, also mindestens gleich 1 ist, so geht aus (15) hervor
, dass, abgesehen von dem Falle $n=1$, $\Delta$ absolut gr\"osser als
1 sein muss, \gesperrt{dass es also ausser dem K\"orper der rationalen
Zahlen keinen K\"orper giebt, dessen Grundzahl $\pm 1$ w\"are.} Dies
ist der von \gesperrt{Minkowski} zuerst erbrachte, fr\"uher lange vergeblich
 gesuchte Beweis dieses wichtigen Satzes\footnote{In einem Briefe an \gesperrt{Hermite} (Comptes rendus d. Pariser Akademie,
26. Januar 1891) und auch in dem oben erw\"ahnten Buche findet \gesperrt{Min-}
\gesperrt{kowski} eine weit engere Grenze, indem er nicht nur das Minimum
der quadratischen Formen, sondern noch anderer Functionen betrachtet,
wie $\eta_1^p + \eta_2^p + \cdots + \eta_n^p$, worin $p$ eine beliebige positive Zahl ist, die
nicht einmal ganz zu sein braucht. F\"ur $p=1$ ergiebt sich, wenn $n-\nu$
die Anzahl der conjugirt imagin\"aren Paare ist:
\[ N_a(\psi) < \left(\frac{4}{\pi}\right)^{n-\nu} \frac{1.2.3\ldots n}{n^n}\, \sqrt{\pm \Delta}. \]
Daraus lassen sich noch weitere Schl\"usse ziehen, wie die: Die Grundzahl
 eines quadratischen K\"orpers muss gr\"osser als 4 oder keiner als % [sic: Weber prints ``keiner'' for ``kleiner''.]
$-2$ sein; die Grundzahl eines cubischen K\"orpers muss gr\"osser als 20
oder kleiner als $-12$ sein etc.}.

\begin{center}
\S.\,157.

\gesperrt{Primfactoren der nat\"urlichen Primzahlen.}
\end{center}

Eine genauere Untersuchung der im \S.\,146 erkl\"arten Basisform
 $\tau$ von $\mathfrak{o}$ des K\"orpers $\Omega$ soll uns nun das Mittel geben,
%CONT>>
